%% ma: L10-B12-0002
%% lop: 10
%% chuong: 5
%% bai: 12
%% dang: 10.12.2
%% loai: TL
%% muc: VD
%% dap_an: "-"
%% nguon: "chuyen-de-cac-so-dac-trung-cua-mau-so-lieu-khong-ghep-nhom-toan-10-chuong-trinh-moi.pdf (GV. Phan Nhật Linh) (P01), Bài 01 Bài tập tự luận 2"
%% kien_thuc: "Lớp 10 Bài 12 (số gần đúng, sai số tuyệt đối, độ chính xác, sai số tương đối, quy tròn số)"
%% trang_thai: chinh-thuc
%% ghi_chu: "Đề gốc thiếu điều kiện x > 0 (lời giải gốc dùng x + 2 > 2), đã bổ sung. Tự duyệt theo yêu cầu giáo viên 22/10/2026."
\begin{ex}
Giả sử $x>0$ là một giá trị gần đúng của $\sqrt5$. Xét số $a=\dfrac{2x+5}{x+2}$. Chứng minh rằng $\left|a-\sqrt5\right|<\left|x-\sqrt5\right|$, tức là nếu lấy $a$ làm giá trị gần đúng của $\sqrt5$ thì được độ chính xác cao hơn so với lấy $x$.
\loigiai{Đặt $u=x-\sqrt5$ (giả sử $u\ne0$). Ta có
\[a-\sqrt5=\dfrac{2x+5-\sqrt5x-2\sqrt5}{x+2}=\dfrac{\left(2-\sqrt5\right)\left(x-\sqrt5\right)}{x+2}=\dfrac{\left(2-\sqrt5\right)u}{x+2}.\]
Do $x>0$ nên $\left|a-\sqrt5\right|=|u|\cdot\dfrac{\sqrt5-2}{x+2}<|u|\cdot\dfrac{\sqrt5-2}{2}<|u|=\left|x-\sqrt5\right|$ (vì $\sqrt5-2<2$).}
\end{ex}
